%% kapa: exam=1 round=25 year=2014 period=1 form=A subject=economics subject_ko=경제원론 seq=27 number=67
%% source: Q-Net 감정평가사 기출문제 (cst00309) artlSeq=5101207 file=extracted/1008/5101207/2014년 제25회 감정평가사 1차시험 기출문제/1교시(A형).hwp
%% conversion: kordoc 4.24.51 -> tools/gen1.py (auto); review=proofread-sonnet; flags=none
\begin{question}{67}
화폐수요함수가 $\frac{ M } { P } =500+0.2Y-1,000i$이다. $Y$=1,000, $i$=0.1일 때, $P$=100과 $P$=200이라면 화폐유통속도는 각각 얼마인가? (단, $M$은 통화량, $P$는 물가, $Y$는 실질국민소득, $i$는 명목이자율, $V_{ 1 }$은 $P$=100일 때 화폐유통속도, $V_{ 2 }$는 $P$=200일 때 화폐유통속도이다.)
\begin{choices}
  \item $V_{ 1 } = \frac{ 5 } { 6 }$, $V_{ 2 } = \frac{ 5 } { 6 }$
  \item $V_{ 1 } = \frac{ 5 } { 6 }$, $V_{ 2 } = \frac{ 10 } { 6 }$
  \item $V_{ 1 } = \frac{ 10 } { 6 }$, $V_{ 2 } = \frac{ 10 } { 6 }$
  \item $V_{ 1 } = \frac{ 10 } { 6 }$, $V_{ 2 } = \frac{ 20 } { 6 }$
  \item $V_{ 1 } = \frac{ 20 } { 6 }$, $V_{ 2 } = \frac{ 20 } { 6 }$
\end{choices}
\end{question}
\answer{3}
